\section{Key Finding: Structure Beats Parameters}
\label{sec:key-finding}

\subsection{The Maximum Partisan Effect is 0.3 Seats}

Across all five parameters and all non-baseline grid points, the maximum
deviation in $D_{\rm nat}$ from the baseline is $\pm 0.2$ seats per parameter.
In the worst case, an adversary simultaneously sets all five parameters
to their most Republican-favorable values:
$\ufactor = 5\%$, $\acounty = 10$, $T = 1{,}000$, $\aswing = 1.20$,
$\wvra = 0.6$.

The combined effect (assuming additive, which overstates the real
correlation) is at most:
\[
  \Delta D_{\rm nat}^{\rm max} = |{-0.2}| + |{-0.1}| + 0 + |{-0.1}| + |{-0.1}|
  = 0.5 \text{ seats}.
\]
In practice, the joint effect is smaller due to correlation and
saturation; empirical combined sweeps yield $\Delta D_{\rm nat} \leq 0.3$.

\subsection{Comparison to Algorithm-Level Effects}

From B.0~\citep{b0paper}, the bakeoff between unweighted bisection and
GeoSection in Wisconsin produced a partisan gap of 4.2 percentage points
in seat share (approximately one seat shift in an 8-seat state).
Scaled to a 435-seat national delegation, this represents approximately
$435 \times 0.042 \approx 18$ seats.

The parameter effect ($\leq 0.3$ seats nationally) is approximately
$60\times$ smaller than the algorithm-structure effect.
This vindicates the B.0 conclusion: algorithm choice is the key lever;
parameter tuning within a fixed algorithm is not.

\subsection{Why Partisanship is Parameter-Insensitive}

The insensitivity of partisan outcomes to parameter tuning reflects two
structural features of the redistricting problem:

\paragraph{Population neutrality.}
Parameters like $\ufactor$ and $T$ affect \emph{how well} the algorithm
satisfies compactness and population constraints, not \emph{which party}
benefits.
Population is not correlated with partisan preference at the tract level
in a way that would make compactness parameters partisan.

\paragraph{Geographic continuity.}
The bisection tree structure is determined by the prime factorisation of $k$
(B.11~\citep{b11paper}) and the geographic graph topology.
Parameter changes alter the leaf-level cut quality within a fixed tree structure,
not the tree structure itself.
The partisan distribution of tracts is not sensitive to leaf-level cut quality.

\paragraph{Marginal seat geography.}
Partisan outcomes depend primarily on how the algorithm handles
geographically marginal regions---the competitive swing districts.
These regions are defined by geography and partisan demography, not by
algorithm parameters.
A parameter change that uniformly shifts district boundaries by 1--2 tracts
will not systematically favor one party in swing regions.

\subsection{The $\ufactor$ Compactness Effect}

Unlike partisanship, compactness \emph{is} parameter-sensitive.
The $\ufactor$ parameter drives $\overline{\PP}$ variation of $\pm 0.010$
(a 3.2\% relative change), with smaller $\ufactor$ producing more compact
districts.

This has a statutory implication: the DIA should specify $\ufactor = 1\%$
(or the smallest value compatible with runtime requirements) as the default.
The choice of $\ufactor$ is consequential for compactness but not for
partisanship.
